\section{Same-depth expansions and the measure comparison}
\label{sec:pairs}

We compare expansions at many nodes of a single generation. Almost
every pair of punctures is sufficiently separated. A level-continuum
argument then forces one of the two charts to spend a definite amount
of oscillation near the other's puncture. Pulling those neighborhoods
back gives more total measure than their overlap permits.

\begin{proposition}
\label{prop:no-subexponential-records}
The moduli $M(n)$ cannot be unbounded while satisfying
\[
 \limsup_{n\to\infty}M(n)^{1/n}\leq1.
\]
\end{proposition}

We prove this proposition through the pair comparison and counting
estimates below. Fix the flow of \Cref{prop:probability-flow}, including $D$, and a
compact neighborhood of its centers in $U$. Use the point $z_0$ chosen
in the flow construction, a positive distance from that neighborhood.
At generation $n$ write
\[
 k=k_0+Dn,\qquad R=a^{-k},\qquad p_i=h_i z_0,
\]
where $h_i=v_i^{-1}$ is the expansion already chosen for node $i$.
The centers $x_i$ are $HR$-separated, and
\begin{equation}
 P_i\le a^{-\gamma Dn}\frac{\Delta_i}{\Delta_0},
 \qquad \frac{P_j}{P_A}\le a^{-\gamma Dm}
 \quad\text{if $j$ lies $m$ generations below $A$}.
 \label{eq:pairs-flow-input}
\end{equation}
We use the arbitrarily large generations having good mass at least
$1/2$, supplied by \Cref{lem:good-levels}.

A good node has a pattern $(\mathcal B^-,\mathcal B^+)$ from a fixed
finite list. Each target ball contains a compact path with endpoint
$u\circ h_i^{-1}$-variation at least $\kappa\Delta_i$, and lies in
the image of its outer source ball. The two target range intervals
have distance at least $\Delta_i/4$. Throughout, a \emph{range interval}
means the interval between the minimum and maximum; metric balls need
not be connected.

All geometric and counting constants below are independent of the
generation, nodes, and approximating index $N$. Constants and thresholds
in convergence statements may depend on a pair after it has been fixed.

\subsection{A probability class and separated punctures}

\begin{lemma}[Bounded puncture multiplicity]
\label{lem:puncture-multiplicity}
For every fixed $C>0$, the balls
$\{\overline B(p_i,CR):i\text{ at generation }n\}$ have uniformly
bounded multiplicity.
\end{lemma}

\begin{proof}
Centers assigned to the same depth-$k$ vertex have mutual distance
$O(R)$. Their $HR$-separation and \Cref{lem:packing} bound their number.
If several puncture balls share a point, their punctures have mutual
products at least $k-O_C(1)$. By \eqref{eq:inverse-prefix} and
\Cref{lem:prefix}, all corresponding inverse vertices lie in one
bounded Cayley-graph ball. There are boundedly many such vertices,
each with the bounded multiplicity just proved.
\end{proof}

There are at most $Ca^{Qk}$ nodes. Discarding those with
$P_i<a^{-2Qk}$ loses probability at most $Ca^{-Qk}$, so the remaining
good nodes have mass at least $1/3$ for large $n$. Partitioning by
pattern and into the $O(k)$ dyadic probability classes gives a class
$\mathcal I$ with
\begin{equation}
 P_*\le P_i\le2P_*\quad(i\in\mathcal I),
 \qquad \sum_{i\in\mathcal I}P_i\ge\frac c k.
 \label{eq:pairs-class}
\end{equation}
Writing $\mathcal N=|\mathcal I|$, we have
\begin{equation}
 \mathcal NP_*\ge\frac{c}{2k}.
 \label{eq:pairs-class-size}
\end{equation}
Also $P_i\le a^{-\gamma Dn}$ by \eqref{eq:pairs-flow-input}, so
$\mathcal N\ge ck^{-1}a^{\gamma Dn}\to\infty$.
We will use punctures separated at a scale $a^{-l}$ much larger than
$R$. The final source-overlap estimate carries a factor $a^{Ql}$,
so we arrange $Ql\leq\gamma Dn$. Put
\[
 \zeta=\gamma/Q,\qquad m=\lfloor\zeta n\rfloor,\qquad l=Dm.
\]
Both $l$ and $k-l$ tend to infinity.

\begin{lemma}[Probability near a puncture]
\label{lem:near-puncture-probability}
For every generation-$n$ node $i$,
\begin{equation}
 \sum_{\substack{j\text{ at generation }n\\d(p_j,p_i)\le a^{-l}}}P_j
 \le Ca^{-\gamma Dm}.
 \label{eq:near-puncture-probability}
\end{equation}
\end{lemma}

\begin{proof}
Fix an ancestor $A$ at generation $n-m$, whose depth is
$k-l=k_0+D(n-m)$. We claim that only boundedly
many descendants $j$ have $d(p_j,p_i)\le a^{-l}$. Split a geodesic word
\[
 v_j^{-1}=w_ju'_j,\qquad |w_j|=l,\qquad |u'_j|=k-l.
\]
The puncture condition and \eqref{eq:inverse-prefix} give
$(p_j|p_i)_e\ge l-C$ and $(v_j^{-1}|p_j)_e\ge k-C$.
Thus $w_j$ lies a bounded distance from the depth-$l$ point on a fixed
ray to $p_i$, leaving boundedly many possibilities.

Reversing and inverting \emph{the same word} gives the geodesic
$v_j=(u'_j)^{-1}w_j^{-1}$. Its depth-$(k-l)$ prefix is therefore
$(u'_j)^{-1}$. Since
\[
 d(x_j,x_A)\le2(B_0+1)r_{k-l},
\]
thinness and \Cref{lem:prefix} put this prefix a bounded distance from
the depth-$(k-l)$ point on a fixed ray to $x_A$. Hence $u'_j$ also has
boundedly many possibilities. Their products determine boundedly
many vertices $v_j$, and the first step of
\Cref{lem:puncture-multiplicity} bounds the nodes assigned to each.
We count group elements here, so nonuniqueness of geodesic spellings
causes no additional multiplicity. The source and inverse prefixes
together account for the full word.

Each such descendant has $P_j\le P_Aa^{-\gamma Dm}$ by
\eqref{eq:pairs-flow-input}. Sum first over these descendants and
then over ancestors, whose probabilities sum to one.
\end{proof}

For each $i\in\mathcal I$, the fraction of its class with puncture
distance at most $a^{-l}$ is bounded by
\[
 \frac{Ca^{-\gamma Dm}}{\mathcal NP_*}
 \le Ck a^{-\gamma Dm}=o(1).
\]
This includes the diagonal. Consequently at least $c\mathcal N^2$
unordered pairs satisfy
\begin{equation}
 d_{ij}:=d(p_i,p_j)>a^{-l}.
 \label{eq:far-pair}
\end{equation}
Call them \emph{far pairs}. Uniformly for these pairs,
$R/d_{ij}\le a^{-(k-l)}\to0$.

\subsection{Comparison on a compact common domain}

Let $K\subset U$ contain all centers, with its closed
$\epsilon$-neighborhood still in $U$. The upper puncture bound gives
\begin{equation}
 d(h_i y,p_i)\ge A_hR/2
 \quad\Longrightarrow\quad
 d(x_i,y)\le D_y\le2C/A_h.
 \label{eq:exterior-preimage}
\end{equation}
Choose $A_h$ so that $2C/A_h<\epsilon/2$ and $A_h/2>2C/d_1$, and set
\[
 E_i=\overline B(p_i,A_hR).
\]
Every $h_iC_{N,t}$ reaches $B(p_i,A_hR/2)$: the diameter bound supplies
a source point at distance at least $d_1/2$ from $x_i$, whose image
has puncture distance at most $2CR/d_1$.

The target balls lie a uniform distance $\delta_*>0$ from $p_i$.
Indeed, their preimages lie in the outer source balls, where
$D_y\le(B_0+2)R$, so the lower puncture bound applies for large $n$.
Because the pattern is shared, this holds for every $i\in\mathcal I$.
For large $n$, uniformly for far pairs,
\[
 2A_hR<a^{-l},\qquad A_hR<\delta_*/2.
\]
Thus $E_i,E_j$ are disjoint, each lies compactly in the other's
transformed domain by \eqref{eq:exterior-preimage}, and the pattern
balls avoid both. These thresholds are independent of $N$.

\begin{lemma}[Pair alternative]
\label{lem:pair-alternative}
Every far pair at a sufficiently large generation has an endpoint
$i$ satisfying
\begin{equation}
 \osc_{h_i^{-1}E_j}u\ge\frac{\kappa}{16}\Delta_i.
 \label{eq:pair-alternative}
\end{equation}
\end{lemma}

\begin{proof}
Write $F_i=u\circ h_i^{-1}$ and $F_j=u\circ h_j^{-1}$ on their open
domains. If $\osc_{E_i}F_j\ge\Delta_j/16$, endpoint $j$ qualifies.
Here we use $\kappa\le1$.
Otherwise, one pattern ball $\mathcal B^s$ has its $F_j$ range interval
at distance at least $\Delta_j/16$ from that on $E_i$. Indeed, if
neither did, the two pattern intervals would have distance less than
$3\Delta_j/16$, contradicting their separation by $\Delta_j/4$.

Node $i$ supplies a compact path $\eta\subset\mathcal B^s$ with
endpoint $F_i$-gap at least $\kappa\Delta_i$. Suppose
$\osc_{E_j}F_i<\kappa\Delta_i/4$. Removing this latter range interval
from the endpoint gap leaves a positive-length closed interval
$I\subset(0,1)$ strictly between the endpoint values and a positive
distance from the range on $E_j$.
To obtain positive margins, choose a smaller closed interval inside
one of the remaining open intervals. Its length need not be uniform
over pairs.

\emph{Fix the pair and $I$ before letting $N\to\infty$.}
Put $F_{i,N}=H_N\circ h_i^{-1}$ and $F_{j,N}=H_N\circ h_j^{-1}$.
The fixed range margins and local uniform convergence imply, for
every $t\in I$ and all sufficiently large $N$, that $h_iC_{N,t}$ hits
$\eta$ and avoids $E_j$. It also reaches the inner $i$-puncture ball.
For the first assertion, the endpoint values still strictly straddle
all of $I$, so \Cref{lem:level-continua} applies to the preimage path.
For the second, the range on $E_j$ remains disjoint from $I$.
By \Cref{lem:stop-continuum}, it has a subcontinuum $K_{N,t}$ containing
a hit on $\eta$, meeting the sphere $d(z,p_i)=A_hR$, and avoiding
$B(p_i,A_hR)$. It remains in the level $F_{i,N}=t$ and avoids $E_j$.
In particular,
\[
 K_{N,t}\subset K_{ij}:=
 Z\setminus\bigl(B(p_i,A_hR)\cup B(p_j,A_hR)\bigr)
 \Subset h_iU\cap h_jU.
\]
Compact containment follows from \eqref{eq:exterior-preimage} with
half the radius. The sphere hit belongs to $E_i$; no identification
of that sphere with the boundary of a metric ball is needed.

Uniform convergence on $\eta\cup E_i$ preserves the $F_j$ range gap,
giving $\osc_{K_{N,t}}F_{j,N}\ge\Delta_j/32$ uniformly for $t\in I$.
The finite-cell comparison, \Cref{lem:cross-energy}, now gives
\[
 \frac{\Delta_j}{32}|I|\le
 \sum_{\substack{q\in T_N\\B(q,2r_N)\cap K_{ij}\ne\varnothing}}
 \osc_{\overline B(q,2r_N)}F_{i,N}\,
 \osc_{\overline B(q,2r_N)}F_{j,N}.
\]
For large $N$ these closed cells lie in the common domain. Both squared
oscillation sums tend to zero by \Cref{lem:local-oscillation}, so
Cauchy--Schwarz contradicts the positive left side. Therefore
$\osc_{E_j}F_i\ge\kappa\Delta_i/4$, proving the alternative.
The argument applies separately to every pair above the geometric
threshold; no convergence rate uniform over pairs is required.
\end{proof}

\subsection{Cost balls and total multiplicity}

We now turn the pair alternative into a measure bound and control the
overlap of the sets being charged. The two indices have distinct roles:
with $i$ fixed, the target punctures control the number of partners $j$;
at a source point, separation of the centers controls the possible
indices $i$. Both estimates will apply to the enlarged balls required
by \eqref{eq:measure}.

For an ordered far pair put $z_{ij}=h_i^{-1}p_j$. The relevant source
scales are
\[
 \begin{array}{c|c}
 \text{puncture denominator }D_{z_{ij}}&\asymp R/d_{ij}\\
 \text{radius of the charged neighborhood}&R^2/d_{ij}^2.
 \end{array}
\]
The radius divided by the denominator is $O(R/d_{ij})\to0$. This separation of scales is what
allows a fixed enlargement of the charged neighborhood without losing
control of its image.

If $y\in h_i^{-1}E_j$, then
$d(h_i y,p_i)\ge d_{ij}-A_hR\ge d_{ij}/2$, whence
$D_y\le CR/d_{ij}$. These points approach $x_i$ uniformly, so the
two-sided puncture bound applies and gives
\begin{equation}
 D_{z_{ij}}\asymp R/d_{ij}.
 \label{eq:cost-denominator}
\end{equation}
The lower expansion bound and $d(h_i y,p_j)\le A_hR$ then give
$d(y,z_{ij})\le CR^2/d_{ij}^2$. Fixing a sufficiently large $C_2$, put
$s_{ij}=C_2R^2/d_{ij}^2$; thus
\begin{equation}
 h_i^{-1}E_j\subseteq\overline B(z_{ij},s_{ij}).
 \label{eq:cost-containment}
\end{equation}
Choose the enlargement constant $C_\mu$ in \eqref{eq:measure}
uniformly on a compact neighborhood of the centers and define
\[
 \mathcal D_{ij}=\overline B(z_{ij},C_\mu s_{ij}).
\]
Their centers stay in this neighborhood and their radii tend uniformly
to zero. If $i$ satisfies \Cref{lem:pair-alternative}, \eqref{eq:measure} gives
\begin{equation}
 \mu(\mathcal D_{ij})\ge c\kappa^2\Delta_i^2.
 \label{eq:cost-lower-bound}
\end{equation}

The enlarged balls still have small images. Indeed, $D_y$ is
$1$-Lipschitz and
\[
 \frac{C_\mu s_{ij}}{D_{z_{ij}}}
 \le C\frac R{d_{ij}}\le Ca^{-(k-l)}\longrightarrow0.
\]
Thus $D_y\asymp D_{z_{ij}}$ on $\mathcal D_{ij}$, and the upper
expansion bound yields
\begin{equation}
 h_i\mathcal D_{ij}\subseteq\overline B(p_j,C_3R).
 \label{eq:cost-image}
\end{equation}
Figure~\ref{fig:pair-cost} shows the charged neighborhood and the image
control that survives the fixed enlargement.

\begin{figure}[htbp]
\centering
\begin{tikzpicture}[x=1cm,y=1cm,font=\small,
  line cap=round,line join=round,>=Stealth]
  \definecolor{costblue}{RGB}{54,101,125}
  \definecolor{costochre}{RGB}{164,111,48}
  \node[font=\bfseries] at (2.75,4.95) {Source};
  \node[font=\bfseries] at (11.40,4.95) {Target};

  \fill[costblue!5] (3.25,2.80) circle (1.10);
  \draw[costblue,thick] (3.25,2.80) circle (1.10);
  \draw[costblue,dashed,thick] (3.25,2.80) circle (.52);
  \node[costblue,anchor=south] at (3.25,4.00) {$\mathcal D_{ij}$};
  \fill[black!70] (.55,2.80) circle (1.8pt);
  \node[anchor=north] at (.55,2.54) {$x_i$};
  \draw[black!40,thin] (3.25,2.73)--(3.25,1.47);
  \fill[costblue] (3.25,2.80) circle (1.8pt);
  \node[anchor=north] at (3.25,1.37) {$z_{ij}$};
  \node[align=center] at (2.75,.57)
    {$h_i^{-1}E_j\subseteq\overline B(z_{ij},s_{ij})$};
  \node[align=center] at (2.75,.05)
    {$s_{ij}\asymp R^2/d_{ij}^{\,2}$};

  \draw[black!85,thick,->] (5.55,2.80)--(7.25,2.80)
    node[midway,above=5pt] {$h_i$};

  \fill[costochre!7] (9.00,2.80) circle (.50);
  \fill[costochre!7] (13.20,2.80) circle (.50);
  \draw[costochre,thick] (9.00,2.80) circle (.50);
  \draw[costochre,thick] (13.20,2.80) circle (.50);
  \draw[costblue,dashed,thick] (13.20,2.80) circle (1.10);
  \node[costochre,anchor=south] at (9.00,3.57) {$E_i$};
  \node[costochre,anchor=west] at (14.50,3.65) {$E_j$};
  \draw[costochre!75,thin] (14.38,3.57)--(13.51,3.19);
  \draw[black!40,thin] (9.00,2.73)--(9.00,1.47);
  \draw[black!40,thin] (13.20,2.73)--(13.20,1.47);
  \fill[costochre] (9.00,2.80) circle (1.8pt);
  \fill[costochre] (13.20,2.80) circle (1.8pt);
  \node[anchor=north] at (9.00,1.37) {$p_i$};
  \node[anchor=north] at (13.20,1.37) {$p_j$};
  \draw[black!85,<->] (9.00,.76)--(13.20,.76)
    node[midway,below=3pt] {$d_{ij}\gg R$};
  \node[costblue,align=center] at (11.10,-.30)
    {$h_i\mathcal D_{ij}\subseteq\overline B(p_j,C_3R)$};
\end{tikzpicture}
\caption{The measure cost of a far pair. The dashed source ball contains
$h_i^{-1}E_j$ and has radius $s_{ij}\asymp R^2/d_{ij}^2$.
Its fixed enlargement $\mathcal D_{ij}$, drawn solid, still maps into
the dashed target ball of radius $C_3R$. Here $p_i=h_i z_0$, $z_{ij}=h_i^{-1}p_j$, and
$E_i,E_j$ are the solid puncture balls of radius $A_hR$.
The circles are schematic; no connectedness or boundary regularity of
metric balls is assumed.}
\label{fig:pair-cost}
\end{figure}

For fixed $i$, \Cref{lem:puncture-multiplicity} therefore bounds the
multiplicity of $\mathcal D_{ij}$ as $j$ varies. Also
\begin{equation}
 \mathcal D_{ij}\subseteq\overline B(x_i,C_4Ra^l),
 \label{eq:cost-source}
\end{equation}
because $d(x_i,z_{ij})\le CR/d_{ij}$ and the cost radius is
$o(R/d_{ij})$. The $HR$-separation of the $x_i$ and
\Cref{lem:packing} bound the multiplicity of these larger balls by
$Ca^{Ql}$.

\begin{proof}[Conclusion of the proof of \Cref{prop:no-subexponential-records}]
Assign each unordered far pair to a qualifying endpoint. At any point,
\eqref{eq:cost-source} allows only $Ca^{Ql}$ paying indices $i$, and
\eqref{eq:cost-image} allows boundedly many partners for each. Integrating
this multiplicity bound gives
\begin{equation}
 \sum_{\{i,j\}\text{ far, assigned to }i}\mu(\mathcal D_{ij})
 \le Ca^{Ql}\mu(Z).
 \label{eq:cost-total-upper}
\end{equation}
This counts the enlarged balls directly and uses no invariance of $\mu$.
The assignment may be arbitrary when both endpoints qualify: all
lower bounds hold for either endpoint, and the multiplicity estimate
already counts every possible ordered contribution.

Conversely, every paying endpoint satisfies
\[
 \Delta_i^2\ge\Delta_0^2a^{2\gamma Dn}P_i^2
 \ge\Delta_0^2a^{2\gamma Dn}P_*^2.
\]
There are at least $c\mathcal N^2$ far pairs, so
\eqref{eq:cost-lower-bound} and \eqref{eq:pairs-class-size} give total cost
\begin{equation}
 c\mathcal N^2\kappa^2\Delta_0^2a^{2\gamma Dn}P_*^2
 \ge c'k^{-2}a^{2\gamma Dn}.
 \label{eq:cost-total-lower}
\end{equation}
Since $\kappa,\Delta_0>0$ are fixed and
$Ql=QD\lfloor\gamma n/Q\rfloor\le\gamma Dn$, comparison with
\eqref{eq:cost-total-upper} yields
\[
 a^{\gamma Dn}\le Ck^2.
\]
This contradicts $k=k_0+Dn$, excluding the unbounded subexponential case.
\end{proof}
